\section{Introduction}
\label{sec:intro}

Individual animal re-identification (Re-ID) enables non-invasive population monitoring, movement analysis, and behavioral research by associating observations of the same animal across sightings~\cite{tuia2022perspectives,vcermak2024wildlifedatasets}. Reliable identification requires preserving subtle individual cues, including coat patterns, scars, and body markings, despite changes in pose, viewpoint, illumination, and time. In multi-species settings, these cues also vary substantially across animal groups. A general-purpose representation must therefore accommodate broad differences between species while retaining the fine-grained distinctions that separate individuals within each species.

Recent work has advanced shared representations for animal Re-ID. WildlifeDatasets introduced MegaDescriptor and a unified framework for learning across heterogeneous wildlife datasets, while OpenAnimals investigated the transfer of person Re-ID techniques to animals~\cite{vcermak2024wildlifedatasets,hou2025openanimals}. However, multi-species training contains a supervisory structure that identity-only objectives leave implicit: individual identities are nested within species. ArcFace~\cite{deng2019arcface} and Dynamic ArcFace~\cite{jiao2021dynarcface} promote identity separation through fixed and adaptive angular margins, respectively, while SphereFace2~\cite{wen2022sphereface2} formulates recognition through one-versus-all binary classification. In their standard identity-only formulations, none explicitly supervises species membership. Different individuals receive distinct identity labels whether they belong to the same species or to different species, leaving the relationship between coarse biological categories and fine-grained identity discrimination unspecified. This motivates an objective that explicitly incorporates both levels of supervision into a shared representation.

The visual evidence used to learn this representation is equally consequential. Backgrounds, capture locations, and acquisition conditions can correlate with identity, allowing contextual information to support retrieval without establishing reliable individual recognition. Rueda-Toicen et al.~\cite{ruedatoicen2026jaguar} demonstrate this concern in jaguar Re-ID through separate diagnostics of background reliance and laterality, showing that retrieval performance alone does not establish which evidence drives a match. In heterogeneous multi-species collections, species labels may also correlate with dataset-specific scenery and imaging conditions. Adding species supervision can therefore leave an opportunity for contextual shortcuts unless the training inputs emphasize the animal itself.

We introduce \textbf{ApexID}, a multi-species animal Re-ID approach centered on a \textbf{hierarchical species--identity loss}. The objective combines individual-level discrimination with explicit species-level supervision in a shared embedding space. Its identity component encourages separation among individuals, including those belonging to the same species, while its species component supplies a complementary coarse-category learning signal. This formulation uses the observed species--identity hierarchy without requiring a complete biological taxonomy. We train ApexID on segmented animal foregrounds to restrict access to surrounding scene content and concentrate learning on animal appearance. The two design choices address complementary aspects of representation learning: the hierarchical objective specifies the label structure used during optimization, and foreground-based training constrains the visual information available to satisfy it.

We evaluate the resulting representation using observation-aware partitions that incorporate temporal and encounter information when available, together with duplicate detection and visual-similarity grouping. Building on prior work on similarity-aware animal Re-ID evaluation~\cite{Adam_2025_CVPR}, these partitions reduce overlap between correlated training and evaluation observations and provide a more demanding test of identity generalization. Our evaluation includes closed-set retrieval and open-world settings involving known-individual recognition, novelty detection, and association of previously unseen individuals. In the causal online setting, discovered identities persist as observations arrive, and each assignment uses only information available at that time. Across the evaluated animal Re-ID tasks, the proposed hierarchical objective improves retrieval performance over ArcFace, Dynamic ArcFace, and SphereFace2 under the reported protocols.

Our contributions are threefold:
\begin{itemize}
    \item We introduce a hierarchical species--identity objective for multi-species animal Re-ID, combining explicit species supervision with fine-grained individual discrimination in a shared representation.
    \item We integrate this objective into ApexID with training on segmented animal foregrounds, pairing hierarchical supervision with an input representation designed to reduce access to background context.
    \item We evaluate ApexID against ArcFace, Dynamic ArcFace, and SphereFace2 using observation-aware partitions, and examine the learned representation through both conventional retrieval and open-world recognition protocols.
\end{itemize}